\section{Demanda Elástica}\label{sec:demanda-elastica}

A elasticidade na demanda por recursos constitui um dos desafios enfrentados pelos provedores de internet. Ao longo do tempo, usuários e aplicações podem necessitar de diferentes quantidades dos recursos oferecidos pela infraestrutura, fazendo com que existam períodos de maior ou menor utilização da rede. Dessa forma, a quantidade de recursos necessária para atender às requisições pode aumentar ou diminuir de acordo com as condições de utilização, caracterizando um cenário de demanda elástica \cite{gomes2020strategies}.

A Figura~\ref{fig:demanda-elastica} apresenta uma representação conceitual do fatiamento dinâmico de redes: serviços com requisitos distintos, como vídeo, jogos e IoT, são atendidos por fatias virtuais cuja demanda de largura de banda varia ao longo do dia. Os valores exibidos têm caráter ilustrativo e não correspondem a medições. A figura evidencia que a quantidade de recursos necessária para atender a um mesmo serviço pode se modificar ao longo do tempo.

\begin{figure}[h!]
    \centering
    \includegraphics[width=\linewidth,keepaspectratio]{figuras/DemandaE.png}
    \caption{Representação do fatiamento de redes e da alocação dinâmica de recursos em ambientes de demanda elástica.}
\label{fig:demanda-elastica}
    \fonte{Elaborado pelo autor com auxílio do ChatGPT (OpenAI, 2026).}
\end{figure}

Diante dessas variações, a alocação dos recursos precisa acompanhar as diferentes condições de demanda apresentadas ao longo do tempo. Nesse contexto, o fatiamento de redes permite que as fatias sejam adaptadas de acordo com as necessidades existentes, possibilitando alterações na utilização dos recursos conforme a demanda aumenta ou diminui \cite{gomes2020strategies}. Dessa forma, a estratégia utilizada para realizar essa adaptação torna-se um fator importante para o gerenciamento dos recursos em ambientes de demanda elástica.
